% Tabla A1-T8. Generada automáticamente; no editar a mano.
% Requiere \usepackage{booktabs} (y graphicx si la tabla es ancha).
\begin{table}[htbp]
    \centering
    \caption{Análisis I: verificación con el radio como variable continua. Respuesta: logaritmo del radio planetario.}
    \label{tab:A1-T8}
    \small
    \resizebox{\textwidth}{!}{%
    \begin{tabular}{llrrrrrrrrrr}
    \toprule
    Modelo & Predictores & $N_{\rm pl}$ & $R^2$ & $R^2_{\rm aj}$ & AIC & BIC & RMSE & MAE & $\Delta R^2$ & $\Delta$AIC & $\Delta$RMSE \\
    \midrule
    C0 & s(logP\_z) & 4\,461 & 0.066 & 0.065 & 3350.9 & 3389.3 & 0.352 & 0.278 & +0.0000 & +0.0 & +0.0000 \\
    C1 & $s($$\log P$$)$ + $M_\star$ + {[Fe/H]} & 4\,461 & 0.227 & 0.226 & 2507.3 & 2558.5 & 0.320 & 0.254 & +0.1617 & -843.6 & -0.0319 \\
    C2 & $s($$\log P$$)$ + $s($$M_\star$$)$ + $s(${[Fe/H]}$)$ & 4\,461 & 0.265 & 0.262 & 2299.8 & 2389.4 & 0.312 & 0.243 & +0.1988 & -1051.1 & -0.0397 \\
    \bottomrule
    \end{tabular}
    }
    \par\smallskip
    {\footnotesize\textit{Nota:} s() denota un spline cúbico. Coeficientes estandarizados en C1 (IC 95 \%, errores agrupados por estrella): st\_mass\_z: 0.123 [0.112, 0.133]; st\_met\_z: 0.055 [0.043, 0.067].}
\end{table}
